\documentclass{book}

\usepackage{amsmath}
\usepackage{amsfonts}
\usepackage{amssymb}
\usepackage{amscd}
\usepackage{amsthm}
\usepackage{graphics}
\usepackage{graphicx}
\usepackage{tikz}
\usepackage{tikz-cd}
\usepackage{bm}

\newtheorem{theorem}{Theorem}
\newtheorem{definition}{Definition}
\newcommand{\bracenom}{\genfrac{\lbrace}{\rbrace}{0pt}{}}

\graphicspath{ {./} }

\title{Topology Munkres - Solutions Manual}
\author{Mingruifu Lin}
\date{September 2023}

\begin{document}

\maketitle

\tableofcontents

\chapter{Set Theory and Logic}

\section{Fundamental Concepts}

\subsection*{Exercise 1}
Okay.

\subsection*{Exercise 2}
(a)
Only $\Rightarrow$ holds.

\noindent
(b)
Only $\Rightarrow$ holds.

\noindent
(c)
True.

\noindent
(d)
Only $\Leftarrow$ holds.

\noindent
(e)
Only $\subseteq$ holds.

\noindent
(f)
Only $\supseteq$ holds.

\noindent
(g)
True:
$$A \cap (B - C)$$
$$= A \cap (B \cap C^c)$$
$$= (A \cap B) \cap C^c$$
$$= \varnothing \cup ((A \cap B) \cap C^c)$$
$$= ((A \cap B) \cap A^c) \cup ((A \cap B) \cap C^c)$$
$$= (A \cap B) \cap (A^c \cup C^c)$$
$$= (A \cap B) \cap (A \cap C)^c$$
$$= (A \cap B) - (A \cap C)$$

\noindent
(h+)
Too lazy.

\subsection*{Exercise 3}
(a)
Contrapositive:

If $x^2 - x \leq 0$, then $x \geq 0$.

\noindent
Converse:

If $x^2 - x > 0$, then $x < 0$.

\noindent
Only the contrapositive is true.

\noindent
(b)
Contrapositive:

If $x^2 - x \leq 0$, then $x \leq 0$.

\noindent
Converse:

If $x^2 - x > 0$, then $x > 0$.

\noindent
None of them are true.

\subsection*{Exercise 4}
(a)
There exists $a \in A$ such that $a^2 \not \in B$.

\noindent
(b)
For all $a \in A$, we have $a^2 \not \in B$.

\noindent
(c)
There exists $a \in A$ such that $a^2 \in B$.

\noindent
(d)
For all $a \not \in A$, we have $a^2 \not \in B$.

\subsection*{Exercise 5}
(a)
Both statement and converse are true by definition.

\noindent
(b)
Statement is false. Converse is true.

\noindent
(c)
Statement is true. Converse is false.

\noindent
(d)
Both statement and converse are true by definition.

\subsection*{Exercise 6}
Too lazy.

\subsection*{Exercise 7}
$$D = A \cap (B \cup C)$$
$$E = (A \cap B) \cup C$$
$$F = A - (B - C)$$

\subsection*{Exercise 8}
If $A$ has $n$ elements, then $\mathcal{P}(A)$ has $2^n$ elements.

\subsection*{Exercise 9}
Done this so many times. In other textbooks.

\subsection*{Exercise 10}
(a)
Yes. $\mathbb{Z} \times \mathbb{R}$

\noindent
(b)
Yes. $\mathbb{R} \times (0, 1]$

\noindent
(c)
No.

\noindent
(d)
Yes. $(\mathbb{R} - \mathbb{Z}) \times \mathbb{Z}$

\noindent
(e)
No.

\section{Functions}

\subsection*{Exercise 1}
(a)
For any function, we have
$$x \in A_0 \quad \Rightarrow f(x) \in f(A_0) \quad \Rightarrow x \in f^{-1}(f(A_0))$$
Now, let $f$ injective. By definition,
$$x \in f^{-1}(f(A_0)) \quad \Rightarrow f(x) \in f(A_0)$$
Suppose for contradiction that $x \not \in A_0$. Denote $y = f(x)$. We know $y$ is an element of $f(A_0)$. In other words, some element of $A_0$ must be able to create it, i.e. $y = f(z)$ for some $z \in A_0$. As we said, since it's not $x$, then there must exist some other $z \in A_0$ where $z \neq x$, such that $f(z) = y$. However, this contradicts the fact that $f$ is injective, i.e. we have $z \neq x \Rightarrow f(z) = f(x)$. Hence, $x \in A_0$ is true.

\noindent
(b)
If $y \in f(f^{-1}(B_0))$, then $y = f(x)$ for some $x \in f^{-1}(B_0)$. By definition, the latter means that $f(x) \in B_0$. We know $y = f(x)$, hence $y \in B_0$.

Now, let $f$ be surjective. If $y \in B_0$, then by surjectivity, we have that $f^{-1}(\{ y \})$ is non-empty. And by definition, for every element $x \in f^{-1}(\{ y \})$, we have $f(x) = y \in B_0$. Hence, $\{ y \} \subseteq f(f^{-1}(\{ y \}))$, thus $B_0 \subseteq f(f^{-1}(B_0))$.

\subsection*{Exercise 2}
(a)
$$x \in f^{-1}(B_0) \quad \Rightarrow f(x) \in B_0 \quad \Rightarrow f(x) \in B_1 \quad \Rightarrow x \in f^{-1}(B_1)$$
(The second "$\Rightarrow$" above uses the assumption given.)

\noindent
(b)
$$x \in f^{-1}(B_0 \cup B_1) \quad \Rightarrow f(x) \in B_0 \cup B_1$$
If $f(x) \in B_0$, then $x \in f^{-1}(B_0)$. Likewise, if $f(x) \in B_1$, then $x \in f^{-1}(B_1)$. Hence, $x \in f^{-1}(B_0) \cup f^{-1}(B_1)$.

For the reverse direction, suppose that
$$x \in f^{-1}(B_0) \cup f^{-1}(B_1)$$
If $x \in f^{-1}(B_0)$, then $x \in f^{-1}(B_0 \cup B_1)$. Likewise, if $x \in f^{-1}(B_1)$, then $x \in f^{-1}(B_0 \cup B_1)$. Hence, $x \in f^{-1}(B_0 \cup B_1)$.

\noindent
(c)
$$x \in f^{-1}(B_0 \cap B_1)$$
$$\Leftrightarrow f(x) \in B_0 \cap B_1$$
$$\Leftrightarrow f(x) \in B_0 \text{ and } f(x) \in B_1$$
$$\Leftrightarrow x \in f^{-1}(B_0) \text{ and } x \in f^{-1}(B_1)$$
$$\Leftrightarrow x \in f^{-1}(B_0) \cap f^{-1}(B_1)$$

\noindent
(d)
$$x \in f^{-1}(B_0 - B_1)$$
$$\Leftrightarrow x \in f^{-1}(B_0 \cap B_1^c)$$
$$\Leftrightarrow f(x) \in B_0 \cap B_1^c$$
$$\Leftrightarrow f(x) \in B_0 \text{ and } f(x) \not \in B_1$$
$$\Leftrightarrow x \in f^{-1}(B_0) \text{ and } x \not \in f^{-1}(B_1)$$
$$\Leftrightarrow x \in f^{-1}(B_0) - f^{-1}(B_1)$$

\noindent
(e)
$$y \in f(A_0) \quad \Rightarrow y = f(x) \mid x \in A_0 \quad \Rightarrow y = f(x) \mid x \in A_1 \quad \Rightarrow y \in f(A_1)$$
(The second "$\Rightarrow$" above uses the assumption given.)

\noindent
(f)
$$y \in f(A_0 \cup A_1) \quad \Rightarrow y = f(x) \mid x \in A_0 \cup A_1$$
If $x \in A_0$, then $y \in f(A_0)$. Otherwise, $y \in f(A_1)$. Hence, $y \in f(A_0) \cup f(A_1)$.

\noindent
(g)
$$y \in f(A_0 \cap A_1)$$
$$\Rightarrow y = f(x) \mid x \in A_0 \cap A_1$$
$$\Rightarrow y = f(x) \mid x \in A_0 \text{ and } y = f(x) \mid x \in A_1$$
$$\Rightarrow y \in f(A_0) \text{ and } y \in f(A_1)$$
$$\Rightarrow y \in f(A_0) \cap f(A_1)$$
Now, suppose that $f$ is injective.
$$y \in f(A_0) \cap f(A_1)$$
$$\Rightarrow y \in f(A_0) \text{ and } y \in f(A_1)$$
$$\Rightarrow y = f(x_0) \mid x_0 \in A_0 \text{ and } y = f(x_1) \mid x_1 \in A_1$$
By injectivity, $x_0 = x_1$, thus
$$\Rightarrow y = f(x) \mid x \in A_0 \text{ and } x \in A_0$$
$$\Rightarrow y = f(x) \mid x \in A_0 \cap A_1$$
$$\Rightarrow y \in f(A_0 \cap A_1)$$

\noindent
(h)
$$y \in f(A_0) - f(A_1)$$
$$\Rightarrow y \in f(A_0) \cap f^c(A_1)$$
$$\Rightarrow y \in f(A_0) \text{ and } y \not \in f(A_1)$$
$$\Rightarrow y = f(x_0) \mid x_0 \in A_0 \text{ and } y \neq f(x_1) \mid \forall x_1 \in A_1$$
Since $y \in \text{range } f|A_0$, there must exist $x \in A_0$ but $x \not \in A_1$ such that $f(x) = y$.
$$\Rightarrow y = f(x) \mid x \in A_0 - A_1$$
$$\Rightarrow y \in f(A_0 - A_1)$$
Now, suppose that $f$ is injective.
$$y \in f(A_0 - A_1)$$
$$\Rightarrow y = f(x) \mid x \in A_0 \cap A_1^c$$
We already know, by conjunction elimination, that $y = f(x) \mid x \in A_0$, hence $y \in f(A_0)$.

Now, for the second part, since $x \in A_1^c$, then $x \not \in A_1$. Suppose for contradiction that $y \in f(A_1)$. Since $y \in f(A_1)$, there must exist $x' \in A_1$ such that $f(x') = y$. We know $x' \in A_1$ and $x \not \in A_1$, hence $x' \neq x$. However, $f(x') = y = f(x)$. This contradicts the injectivity of $f$. Hence, it must be that $y \not \in f(A_1)$.

Thus, $y \in f(A_0)$ and $y \not \in f(A_1)$. Hence,
$$\Rightarrow y \in f(A_0) - f(A_1)$$

If you didn't understand, read the following. The crux to understanding the second part is to realize that $y$ is a value in the codomain. Many $x$ can produce this single $y$. We need to verify that every single $x \in f^{-1}(\{ y \})$ is not contained in $A_1$. If even a single one is found to be in $A_1$, then the entire $y$ is contained in $f(A_1)$. That's where injectivity comes in, guaranteeing that the preimage is a singleton, so that, as soon as we find a particular $x \not \in A_1$ such that $f(x) = y$, then there is no more $x$ to verify.

\subsection*{Exercise 3}
Hell nah.

\subsection*{Exercise 4}
(a)
$$(g \circ f)^{-1}(C_0) = \{ a \in A \mid (g \circ f)(a) \in C_0 \}$$
$$= \{ a \in A \mid \exists b \in B : f(a) = b \text{ and } g(b) \in C_0 \}$$
$$= \{ a \in A \mid \exists b \in B : f(a) = b \text{ and } b \in g^{-1}(C_0) \}$$
$$= \{ a \in A \mid f(a) \in g^{-1}(C_0) \}$$
$$= f^{-1}(g^{-1}(C_0))$$

\noindent
(b)
$$x \neq x' \quad \Rightarrow f(x) \neq f(x') \quad \Rightarrow g(f(x)) \neq g(f(x'))$$

\noindent
(c)
Recall that Munkres uses the word "range" to mean "codomain". Hence, we only require $f$ to be injective. $g$ must be injective on $f(A) \subseteq B$. But it can behave freely elsewhere, i.e. on $B - f(A)$, because these values are never reached by $f$.

\noindent
(d)
For every $c \in C$, we can find $b \in B$ such that $g(b) = c$. And for every $b \in B$, we can find $a \in A$ such that $f(a) = b$. Thus, for every $c \in C$, we can find $a \in A$ such that $c = g(f(a))$. Hence, $g \circ f$ is surjective.

\noindent
(e)
$g$ must be surjective. However, $f$ need not be surjective, because it's possible that $C = g(B_0)$ for some subset $B_0 \subset B$. In other words, $f$ only needs to reach $B_0$ for the entire composite to be surjective.

\noindent
(f)
Too lazy.

\subsection*{Exercise 5}
(a)
Let $f : A \rightarrow B$ have a left inverse $g : B \rightarrow A$ where $g \circ f = i_A$. Suppose for contradiction that $f$ is not injective, i.e. there exists $x \neq x'$ such that $f(x) = f(x')$. Then,
$$g(f(x)) = g(f(x'))$$
But $g \circ f$ is the identity function, hence
$$\Rightarrow x = x'$$
contradicting our assumption. Hence, $f$ is injective.

Now, let $f : A \rightarrow B$ have a right inverse $h : B \rightarrow A$ where $f \circ h = i_B$. For every $b \in B$, we can compute $h(b) \in A$. The identity function tells us that $f(h(b)) = b$. In other words, for every $b \in B$, we found $a \in A$ such that $f(a) = b$, where $a = h(b)$. This is precisely the requirement of surjectivity of $f$.

\noindent
(b, c)
It's really just playing with the codomain.

\noindent
(d)
Yes. Again, when $f$ is not surjective, $f$ need not map to the entirety of $B$, so $g$ is free to do whatever it likes on the remaining portions. And when $f$ is not injective, it can have multiple right inverses. For example, if $x, x' \in A$ both map to $y \in B$, then $h(b)$ can choose between them.

\noindent
(e)
Well, since $f$ is both injective and surjective, then it's bijective.

First, $h$ is surjective. Suppose for contradiction that it's not surjective. Then take any $x \in A - h(B)$, which is non-empty by non-surjectivity of $h$. Every $y \in B$ has already been taken by some $x' \in h(B)$, because $f(h(B)) = B$ by property of identity. Since $x$ is in the domain of $f$, then it must produce something under $f$, say $y \in B$. But then $f(x) = y = f(x')$. We know $x' \neq x$, because they belong in complementary subsets of $A$. This would contradict the injectivity of $f$, hence $h$ is surjective.

Now we can complete the proof. Given $y \in B$, using the surjectivity of $f$, there exists $x \in A$ such that
$$y = f(x)$$
Hence,
$$g(y) = g(f(x)) = x$$
by property of identity. Now, since $h$ is surjective, then there must exist some $y' \in B$ such that
$$h(y') = x$$
It follows that
$$f(h(y')) = f(x)$$
$$\Rightarrow y' = y$$
Hence, $h(y) = h(y') = x = g(y)$ for all $y \in B$, and since they share this same domain, then they are identical functions.

By injectivity of $f$, the $x$ that gives $y = f(x)$ is unique. Hence, $g = h = f^{-1}$ is the inverse function, since it evaluates to this unique $x$.

\subsection*{Exercise 6}
Too lazy.

\section{Relations}

\subsection*{Exercise 1}
Equality is reflexive, symmetric, and transitive. The equivalence classes can be seen if we fix $y - x^2 = c$. Which becomes $y = x^2 + c$. These are parabolas. The equivalence classes differ by the amount shifted vertically.

\subsection*{Exercise 2}
We are basically removing elements from an already related set, and we have to prove that the remaining connections still form an equivalence relation.

For every element $a \in A$, we have $(a, a) \in C$. If $a \in A_0$, then $(a, a) \in A_0 \times A_0$, hence $(a, a) \in C \cap (A_0 \times A_0)$.

Suppose $(a, b) \in C \cap (A_0 \times A_0)$. Clearly, $(b, a) \in C$. And since $a \in A_0$ and $b \in A_0$, then $(b, a) \in A_0 \times A_0$. Thus, $(b, a) \in C \cap (A_0 \times A_0)$.

Suppose $(a, b), (b, c) \in C \cap (A_0 \times A_0)$. Clearly, $(a, c) \in C$. And since $a \in A_0$ and $c \in A_0$, then $(a, c) \in A_0$. Thus, $(a, c) \in C \cap (A_0 \times A_0)$.

\subsection*{Exercise 3}
These are conditional statements. If no pair $(a, b)$ exist, then we can't proceed. In other words, none of these statements enforce existence.

\subsection*{Exercise 4}
(a)
Apply the properties of equality.

\noindent
(b)
Clearly intuitive, but I'm confusing myself with proofs.

\subsection*{Exercise 5}
(a)
$x - x = 0 \in \mathbb{Z}$, so it's reflexive. If $x - y = n \in \mathbb{Z}$, then $y - x = -n \in \mathbb{Z}$, so it's symmetric. If $x - y = n \in \mathbb{Z}$ and $y - z = m \in \mathbb{Z}$, then $x - z = x + (y - y) - z = (x - y) + (y - z) = n + m \in \mathbb{Z}$, so it's transitive.

Obviously, it's a superset of $S$, because
$$(x, y) \in S \quad \Rightarrow y = x + 1 \quad \Rightarrow y - x = 1 \in \mathbb{Z} \quad \Rightarrow (x, y) \in S'$$

Numbers belonging to a same equivalence classes have the same mod 1.

\noindent
(b)
If $a \in A$, then $(a, a)$ is contained in every equivalence relation (because that's a property of equivalence relations). Hence, it's in their intersection.

If $(a, b)$ is contained in the intersection, then it must be contained in each equivalence relation. And each equivalence relation must contain $(b, a)$, hence their intersection contains $(b, a)$ as well.

If $(a, b)$ and $(b, c)$ is contained in the intersection, then they must be contained in each equivalence relation. And each equivalence relation must contain $(a, c)$, hence their intersection contains $(a, c)$ as well.

\noindent
(c)
$(x, x) \in T$ for every $x \in \mathbb{R}$. By the way, $S$ is not an equivalence relation, because it's not symmetric, so since $(x, y) \in S$, then $(x, y) \in T$, and since $T$ is an equivalence relation, then $(y, x) \in T$. In other words, we add the missing symmetry. Transitivity doesn't need to be dealt, because the range $0 < x < 2$ is too small for transitivity.

$T$ basically looks like the real number line, but we glue the interval $(0, 1)$ with $(1, 2)$ sort of in a superposition fashion. There's a hole in the middle, with only $(1, 1) \in T$ (not interval, here I'm expressing a pair).

\subsection*{Exercise 6}
These are parabolas. Higher parabolas are greater than lower parabolas, and within the same parabola, right points are greater than left points.

First, obviously, every pair of points are comparable, by definition of the relation. Second, if two points are identical, then $y_0 - x_0^2 \neq y_1 - x_1^2$ is not satisfied, then $x_0 \neq x_1$ is not satisfied, hence identical points are not comparable.

Third, transitivity is obvious, but tedious. If $(x_0, y_0) < (x_1, y_1) < (x_2, y_2)$, then the second point either sits on the same or higher parabola than the first, and the third either sits on the same or higher parabola than the second. In the same-same case, we have $y_2 - x_2^2 = y_1 - x_1^2 = y_0 - x_0^2$, so we skip the parabola inequality, directly resulting in $x_0 < x_1 < x_2$, hence $x_0 < x_2$. In the same-higher case, we get $y_2 - x_2^2 > y_1 - x_1^2 = y_0 - x_0^2$, hence $y_2 - x_2^2 > y_0 - x_0^2$. In the higher-same case, we get $y_2 - x_2^2 = y_1 - x_1^2 > y_0 - x_0^2$, hence $y_2 - x_2^2 > y_0 - x_0^2$. In the higher-higher case, we simply have $y_2 - x_2^2 > y_1 - x_1^2 > y_0 - x_0^2$, hence $y_2 - x_2^2 > y_0 - x_0^2$.

\subsection*{Exercise 7}
I'm assuming that the restriction excludes points in $A$ outside the restriction:
$$A \rightarrow A \cap B$$
$$C \rightarrow C \cap (B \times B)$$
For comparability:
$$x, y \in A \cap B$$
$$\Rightarrow x, y \in A \text{ and } x, y \in B$$
Since $C$ is an order relation on $A$, then we have $(x, y) \in C$ already. And since $x, y \in B$, then $(x, y) \in B \times B$. Hence,
$$\Rightarrow (x, y) \in C \text{ and } (x, y) \in B \times B$$
$$\Rightarrow (x, y) \in C \cap (B \times B)$$
For nonreflexivity, since $(x, x) \not \in C$, then $(x, x) \not \in C \cap (B \times B)$ even less. Finally, for transitivity
$$(x, y), (y, z) \in C \cap (B \times B)$$
$$\Rightarrow (x, y), (y, z) \in C \text{ and } (x, y), (y, z) \in B \times B$$
$$\Rightarrow (x, z) \in C \text{ and } x, y, z \in B$$
$$\Rightarrow (x, z) \in C \text{ and } (x, z) \in B \times B$$
$$\Rightarrow (x, z) \in C \cap (B \times B)$$

\subsection*{Exercise 8}
Too lazy.

\subsection*{Exercise 9}
The procedure is similar to Exercise 6.

\subsection*{Exercise 10}
(a)
Suppose $0 < x < y < 1$. Then $x^2 < y^2$. Then $1 - x^2 > 1 - y^2$. Then $\frac{1}{1 - x^2} < \frac{1}{1 - y^2}$. Hence, $\frac{x}{1 - x^2} < \frac{y}{1 - y^2}$. If both are negative, i.e. $-1 < x < y < 0$, then $x^2 > y^2$, then $1 - x^2 < 1 - y^2$, then $\frac{1}{1 - x^2} > \frac{1}{1 - y^2}$ since both are negative. Hence, $x$ is scaled by a bigger factor, making it more negative $\frac{x}{1 - x^2} < \frac{y}{1 - y^2}$. Finally, if $x < 0 < y$, then since the denominators are always positive, they maintain their sign, so $\frac{x}{1 - x^2} < 0 < \frac{y}{1 - y^2}$.

\noindent
(b)
Too lazy.

\subsection*{Exercise 11}
Suppose that $b, c$ are both immediate successors of $a$. Then we know $a < b$ and $a < c$. If $b \neq c$, then either $b < c$ or $c < b$ (comparability property). If $b < c$, then $a < b < c$, which violates the fact that $c$ is immediate successor of $a$. The other case is similar. Hence, $b$ and $c$ must be equal. Thus, immediate successors are unique. The procedure is similar for immediate predecessors.

Suppose that $a, b$ are both largest elements. Then $a \leq b$, and $b \leq a$, hence $a = b$. Hence, largest elements are unique. The procedure is similar for smallest elements.

\subsection*{Exercise 12}
(i) Every bottom element of each column has no immediate predecessor. The bottom-left (origin) is the smallest element.

\noindent
(ii) This is basically shifting the line $y = x$, where shifting bottom-right is larger and shifting top-left is smaller. Within a line, top-right is larger. Every point on the bottom or left has no immediate predecessor. There is no smallest element.

\noindent
(iii) This is basically shifting the line $y = -x$, where shifting bottom-left is smaller and shifting top-right is larger. Within a line, top-left is larger. The bottom-left (origin) is the smallest element. Except for the origin, every element has an immediate precedessor.

\subsection*{Exercise 13}
This is a standard exercise in Real Analysis. We'll abbreviate with LUB and GLB. We simply set the GLB of a subset to the LUB of its set of lower bounds. It can be proven that they are equivalent, hence GLB exists by extensionality.

\subsection*{Exercise 14}
(a)
Well, if $C$ is symmetric, then $(a, b) \in C$ implies $(b, a) \in C$. By definition, $D$ contains $(b, a)$ as well, hence every element of $C$ is found in $D$. Likewise, if $(c, d) \in D$, then $(d, c) \in C$. By symmetry, $(c, d) \in C$, hence every element of $D$ is also included $C$. Hence $C = D$.

For the other direction, suppose $C = D$. If $(a, b) \in C$, then by definition of $D$, we have $(b, a) \in D$. Since $C = D$, then $(b, a)$ must be included in $C$ as well. Hence $(a, b) \in C$ implies $(b, a) \in C$, hence $C$ is symmetric.

\noindent
(b)
Given $a, b \in A$, obviously, either $(a, b) \in D$ or $(b, a) \in D$, since either must be in $C$. Also, $(a, a) \not \in D$ for every $a$, because it would imply $(a, a) \in C$, which contradicts Property 2 of the order relation $C$. Finally, if $(a, b), (b, c) \in D$, then $(b, a), (c, b) \in C$. By transitivity of $C$, it follows $(c, a) \in C$, hence $(a, c) \in D$. Since $D$ is comparable, nonreflexive, and transitive, it follows that $D$ is an order relation.

\noindent
(c)
Too lazy.

\subsection*{Exercise 15}
Too lazy.

\section{The Integers and the Real Numbers}
Too lazy. These exercises are closer to Abstract Algebra.

\section{Cartesian Product}

\subsection*{Exercise 1}
Well... assign $(a, b) \rightarrow (b, a)$. Clearly, every $(b, a)$ has a corresponding $(a, b)$, so it's surjective. Also, if $(a, b) \neq (c, d)$, then either $a \neq c$ or $b \neq d$, hence $(b, a) \neq (d, c)$, because otherwise, it would imply $b = d$ and $a = c$ which contradicts the previous, hence it's injective. Thus, the whole thing is bijective.

\subsection*{Exercise 2}
Too lazy.

\subsection*{Exercise 3}
(a)
If $b_1 \times b_2 \times \cdots \in B$, then $b_1 \in B_1, b_2 \in B_2, \ldots$, hence $b_1 \in A_1, b_2 \in A_2, \ldots$, hence $b_1 \times b_2 \times \cdots \in A$.

\noindent
(b)
If $b_0 \in B_i$, consider some $b \in B$ for which its $i$-th coordinate is $b_0$. It must exist, because $B$ is non-empty, hence each $B_i$ is non-empty, hence we can choose fillers for non-$i$ coordinates. By definition of subset, $b$ is matched with some $a \in A$. Hence, $b_0 \in A_i$.

\noindent
(c)
We basically making the crucial step in (b) clear. It's a simple contradiction proof. If one of $A_i$ is empty, then we cannot produce any Cartesian product, contradicting the fact $A$ is non-empty. Hence, each $A_i$ is non-empty.

For the converse, I searched online. We need the Axiom of Choice. It makes sense. Because we want to create a Cartesian product by selecting one element from each $A_\alpha$.

\noindent
(d)
$A \cup B$ is a subset of the Cartesian product of $A_i \cup B_i$. On the other hand, $A \cap B$ is equal to the Cartesian product of $A_i \cap B_i$.

\subsection*{Exercises 4 - 5}
Too lazy.

\chapter{Topological Spaces and Continuous Functions}

\section{Basis for a Topology}

\subsection*{Exercise 1}
$A$ is equal to the union of $U$'s. We know that topol. spaces are closed under arbitrary union.

\subsection*{Exercise 2}
Too lazy.

\subsection*{Exercise 3}
Empty set is included, because $X - \varnothing = X$ which is all of $X$. The set $X$ itself is included, because $X - X = \varnothing$ is countable (it's finite).

Given a collection $W$ of open sets $U$, which satisfy $X - U$ is countable, we check:
$$X - \bigcup_{U \in W} U = \bigcap_{U \in W} (X - U)$$
We know each $X - U$ is countable. (Of course, one of them could satisfy $X - U = X$, but then it would have no effect on the intersection.) Obviously, arbitrary intersection of countable sets is countable, hence the complement of the union is finite, hence the union is inside $\mathcal{T}_c$.

Finally, if $X - U_0$ and $X - U_1$ are countable, then $X - U_0 \cap U_1 = (X - U_0) \cup (X - U_1)$ is countable, hence $U_0 \cap U_1 \in \mathcal{T}_c$. If one of them satisfies $X - U = X$, then the result is all of $X$, so it's still included in $\mathcal{T}_c$. Thus $\mathcal{T}_c$ is a topology.

\subsection*{Exercise 4}
(a)
$\varnothing, X \in \mathcal{T}_\alpha$ for every $\mathcal{T}_\alpha$ already, so they are also included in their intersection. Given a collection of sets in $\bigcap \{ \mathcal{T}_\alpha \}$, each set is also individually inside each $\mathcal{T}_\alpha$. Since the latter contain their union, then their intersection also contains their union. Similarly, given a finite amount of sets in $\bigcap \{ \mathcal{T}_\alpha \}$, each set is individually in each $\mathcal{T}_\alpha$.

The union $\bigcup \{ \mathcal{T}_\alpha \}$ is generally not a topology though. We can take one open set from some $\mathcal{T}_\alpha$ and another from some other $\mathcal{T}_\beta$, and the union of these open sets would not necessarily be included in either topology, hence would not be included in the union of the two topologies.

\noindent
(b)
The smallest surrounding topology must contain every arbitrary union and finite intersection of open sets from each individual topology. It's the smallest, because it satisfies every condition minimally. The largest contained topology is the intersection between the topologies. It's the largest, because it's obvious.

By construction, it's obvious that they are both unique.

\noindent
(c)
Too lazy.

\subsection*{Exercise 5}
The topology generated by $\mathcal{A}$ is formed by taking arbitrary unions of open sets in $\mathcal{A}$. The intersection must contain this topology. Every topology containing $\mathcal{A}$ must contain this topology, because it's the smallest, satisfying the conditions of a topological space minimally. Hence the intersection equals this topology. The same goes if $\mathcal{A}$ is a subbasis, where we take finite intersections alongside arbitrary unions.

\subsection*{Exercise 6 - 7}
Too lazy.

\subsection*{Exercise 8}
(a)
Every open set of the standard topology is the arbitrary union of open intervals. Every point $x$ inside an open interval has a neighborhood $(x - \delta, x + \delta)$ fully contained in the interval. We can find a rational between $x$ and $x - \delta$, as well as another one between $x$ and $x + \delta$. Hence, using these two rationals as bounds for a new interval, it shows that for every point $x$ inside an open set, we can find an interval that contains $x$ which is inside the open set.

\noindent
(b)
The open set $[\sqrt{2}, 2) \in \mathbb{R}_l$ cannot be produced by rational-bounded half-open sets.

\section{The Subspace Topology}

\subsection*{Exercise 1}
Every open set of $A$ equals $A \cap U$ for some open $U \in Y$. And every open $U \in Y$ equals $Y \cap V$ for some open $V \in X$. Hence, every open set of $A$ equals $A \cap (Y \cap V)$ for some open $V \in X$. But since $A \subseteq Y$, then $A \cap Y = A$, so every open set is simply expressed as $A \cap V$ for some open $V \in X$. Hence $A$ has the subspace topology from $X$.

\subsection*{Exercise 2}
The subspace topology from $\mathcal{T}'$ is finer than that of $\mathcal{T}$, but not necessarily strictly finer.

As an example, let the topology $\mathcal{T}'$ be the topology on $\mathbb{R}$ generated by the basis:
$$\mathcal{B} = \{ (a, b) : a, b \in \mathbb{R} \} \cup \{ [0, c) : c \in \mathbb{R} \}$$
which is basically the standard topology but we include a half-open ray starting at zero. It is strictly finer than the standard topology, which is $\mathcal{T}$ in our example. However, taking $Y = [0, \infty)$ as the subset produces two equivalent subspace topologies.

\subsection*{Exercise 3}
Too lazy.

\subsection*{Exercise 4}
Assuming the product topology on $X \times Y$, then every open set in $X \times Y$ is the arbitrary union of a bunch of $U \times V$ for $U, V$ open in $X, Y$ respectively. We know that
$$f \bigcup (U \times V) = \bigcup f(U \times V)$$
Since $\pi_1(U \times V) = U$ which is open, then the union is also open. Hence, $\pi_1$ is open. The procedure is similar for $\pi_2$.

\subsection*{Exercise 5}
(a)
If I understand the problem correctly, $X = X'$ denote the same set, but with different topologies $\mathcal{T}$ and $\mathcal{T}'$. Likewise for $Y = Y'$.

The product topology on $X \times Y$ has basis elements of the form $T \times U$ where $T \in \mathcal{T}$ and $U \in \mathcal{U}$. Since $\mathcal{T} \subseteq \mathcal{T}'$ and $\mathcal{U} \subseteq \mathcal{U}'$, then $T \times U \in \mathcal{T}' \times \mathcal{U}'$. Since every basis of $X \times Y$ is included in $X' \times Y'$, then the latter is finer.

\noindent
(b)
Suppose $T \times U \in X \times Y$ implies $T \times U \in X' \times Y'$. Now, given $T \in \mathcal{T}$, since $Y$ is non-empty, then $T \times Y$ is open in $X \times Y$. Hence, $T \times Y \in X' \times Y'$ open as well. Its projection $T$ is open in $\mathcal{T}'$, hence $\mathcal{T} \subseteq \mathcal{T}'$. The procedure is similar for $\mathcal{U}$ and $\mathcal{U}'$.

\subsection*{Exercise 6}
Well, for every point $x$ inside rational rectangle $(a, b) \times (c, d)$, we can find a subset irrational rectangle $(e, f) \times (g, h)$ containing $x$. The converse is true as well. Hence their topologies are equivalent.

\subsection*{Exercise 7}
No. For example, take a line with a hole, and let the subset be the entire portion after the hole. This subset is convex, but you cannot express its left endpoint (because it's literally the hole). I got this solution from online.

\subsection*{Exercise 8}
It depends on the direction of the line. Too lazy to try all.

\subsection*{Exercise 9}
The dictionary order topology's basis are intervals. Given endpoints $\langle a, b \rangle < \langle c, d \rangle$, any interval in the dictionary topology can be decomposed into
$$\{ a \} \times (b, \infty)$$
$$(a, c) \times \mathbb{R}$$
$$\{ c \} \times (-\infty, d)$$
which belong to $\mathbb{R}_d \times \mathbb{R}$. Conversely, $U \times V$ are the basis for the product topology, where $U \subseteq \mathbb{R}$ are arbitrary (since it's discrete topology) and $V = (p, q) \subseteq \mathbb{R}$ are open intervals. We can express them using the dictionary basis as follows:
$$\bigcup_{x \in U} (\langle x, p \rangle, \langle x, q \rangle)$$
Basically, we are using a bunch of vertical lines. Hence, the two topologies are equal.

Since discrete topology is strictly finer than the standard topology on $\mathbb{R}$, then the product topology must be at least finer as well. In fact, it's also strictly finer, since lines are invalid in the standard topology on $\mathbb{R}^2$, but valid in the other.

\subsection*{Exercise 10}
Too lazy.

\section{Closed Sets and Limit Points}

\subsection*{Exercise 1}
First, $\varnothing = X - X$ and $X = X - \varnothing$, so both are included. Second, suppose $X - C_\alpha$ are open for (possibly uncountable) $\alpha$, then $\bigcup_\alpha (X - C_\alpha) = X - \bigcap_\alpha C_\alpha$, the subtracted part which is an arbitary intersection of elements in $\mathcal{C}$, hence also in $\mathcal{C}$, hence the original thing is in $\mathcal{T}$. Finally, suppose $X - C_1$ and $X - C_2$ are open, then $(X - C_1) \cap (X - C_2) = X - (C_1 \cup C_2)$, the subtracted which is a finite union of elements in $\mathcal{C}$, hence also in $\mathcal{C}$, hence the original thing is in $\mathcal{T}$. Hence, it's a topology.

\subsection*{Exercise 2}
Since $A$ is closed in $Y$, and $Y$ is closed in $X$, then $A = C \cap Y$ for some $C$ closed in $X$. Intersection of closed sets in $X$ is closed in $X$, hence $A$ is closed in $X$.

\subsection*{Exercise 3}
Assuming the product topology. We can decompose the complement as
$$X \times Y - A \times B = ((X - A) \times Y) \cup ((X) \times (Y - B))$$
Since $X - A$ is open, and $Y$ is open, and $X$ is open, and $Y - B$ is open, then the Cartesian product and union are open are well. Since the complement of $A \times B$ is open, then it must be closed itself.

\subsection*{Exercise 4}
Too lazy.

\subsection*{Exercise 5}
Given $x \in X$ limit point of $(a, b)$, then it cannot be smaller than $a$. Otherwise, the ray $(-\infty, a)$ is a neighborhood disjoint from $(a, b)$ which contains $x$, contradicting the fact that $x$ is a limit point. Likewise, $x$ cannot be greater than $b$. Hence, either $x = a$ or $x = b$ or $x \in (a, b)$. This is equivalent to $x \in [a, b]$.

Equality holds when $a$ and $b$ are not isolated from $(a, b)$, which means $a$ and $b$ are in the closure. For example, consider $X = \{ 0 \} \cup (1, \infty)$. Consider the interval $(0, 2)$. It is equal to $(1, 2)$, hence its closure is $(1, 2]$. Notice that $0$ is not included, because it's isolated.

\subsection*{Exercise 6}
(a)
If $a$ is a limit point of $A$, then every neighborhood $U_a$ intersects $A$ at some point. Hence, every $U_x$ intersects $B$ at those points also. Hence, $a$ is a limit poinf of $B$.

\noindent
(b)
$$x \in \overline{A \cup B}$$
$$\Leftrightarrow U_x \cap (A \cup B) \neq \varnothing$$
$$\Leftrightarrow (U_x \cap A) \cup (U_x \cap B) \neq \varnothing$$
$$\Leftrightarrow U_x \cap A \neq \varnothing \text{ or } U_x \cap B \neq \varnothing$$
$$\Leftrightarrow x \in \bar{A} \text{ or } x \in \bar{B}$$
$$\Leftrightarrow x \in \bar{A} \cup \bar{B}$$

\noindent
(c)
$$x \in \bigcup \bar{A}_\alpha$$
$$\Rightarrow x \in \bar{A}_\alpha \text{ for some } \alpha$$
$$\Rightarrow U_x \cap A_\alpha \neq \varnothing$$
$$\Rightarrow U_x \cap \bigcup A_\alpha \neq \varnothing$$
$$\Rightarrow x \in \overline{\bigcup A_\alpha}$$
See next exercise for counter-example.

\subsection*{Exercise 7}
It's false that if every neighborhood intersects $\bigcup A_\alpha$, then one particular $A_\alpha$ is always intersected. Consider $\bigcup (\frac{1}{\alpha}, 1) = (0, 1)$. Obviously, $0$ is in the closure of $(0, 1)$. However, none of the $(\frac{1}{\alpha}, 1)$ is intersected by every neighborhood of $0$, because we simply choose the neighborhood $(-1, \frac{1}{\alpha})$, hence $0$ belongs to none of the individual closures.

\subsection*{Exercises 8 - 9}
Too lazy.

\subsection*{Exercise 10}
For every two distinct points $a$ and $b$, there are two cases. If $b$ is the immediate successor of $a$, then use the neighborhoods $(-\infty, b)$ and $(a, \infty)$. Their intersection equals $(a, b)$ which is empty since they are in successor relationship. In the other case, there must be an element $c$ between $a$ and $b$. Simply use the neighborhoods $(-\infty, c)$ and $(c, \infty)$, which are obviously disjoint. Hence, order topologies are Hausdorff.

\subsection*{Exercise 11}
Suppose $\langle a, b \rangle, \langle c, d \rangle \in X \times Y$. Since $X$ and $Y$ are Hausdorff, then we can find disjoint neighborhoods $U_a$ and $U_c$, as well as disjoint neighborhoods $U_b$ and $U_d$. Then $U_a \times U_b$ and $U_c \times U_d$ are disjoint in $X \times Y$ but each contain their respective point. Hence, the product is Hausdorff.

\subsection*{Exercise 12}
Let $Y$ be a subspace. If neighborhoods $U_a$ and $U_b$ are disjoint, then neigborhoods $U_a \cap Y$ and $U_b \cap Y$ are disjoint as well.

\subsection*{Exercise 13}
Suppose $X$ is Hausdorff. Suppose for contradiction that $\Delta$ is not closed, i.e. there exists $p = \langle p_1, p_2 \rangle$ that is limit point of $\Delta$, but $p \not \in \Delta$. For every pair of neighborhood $U_{p_1}, U_{p_2} \subseteq X$, their product $U_{p_1} \times U_{p_2} = U_p$ is a neighborhood of $p$, which intersects $\Delta$ since it's a limit point. Hence, some point $a \times a$ is contained in $U_p$. Hence, $a$ is contained in each individual $U_{p_1}$ and $U_{p_2}$, hence their intersection is non-empty. Reminder: this works for every pair of neighborhood of $p_1$ and $p_2$. Thus, this contradicts the fact that $X$ is Hausdorff. Hence, $\Delta$ is closed.

Too lazy for other direction.

\subsection*{Exercise 14}
It converges to every point in $\mathbb{R}$. For example, take any $x \in \mathbb{R}$. Take any neighborhood of $x$. It can only exclude a finite amount of points, so there will always be infinite points of $x_n$ still inside the neighborhood.

\subsection*{Exercise 15}
Since points are closed, then the second point (which is distinct) is not a limit point of the first. Hence, it has a neighborhood not containing the first.

\subsection*{Exercises 16 - 17}
Too lazy.

\subsection*{Exercise 18}
I forgot what an ordered square is.

\subsection*{Exercise 19}
(a)
Suppose for contradiction that there exists a point $x \in \text{Int } A \cap \overline{X - A}$. Then $x$ must be part of an open set contained in $A$, so we can take an open neighborhood $B_x \subseteq A$. Also, $x$ is in the closure of the outside, so by definition every neighborhood intersects the outside: $B_x \cap (X - A) \neq \varnothing$. The latter means that there exists some $p \in B_x$ such that $p \not \in A$, but this contradicts $B_x \subseteq A$, hence there exists no such point $x$. Hence $\text{Int } A$ and $\overline{X - A}$ are disjoint.

For the second part,
$$\text{Int } A \cup \text{Bd } A$$
$$= \text{Int } A \cup (\bar{A} \cap \overline{X - A})$$
$$= (\text{Int } A \cup \bar{A}) \cap (\text{Int } A \cup \overline{X - A})$$
$$= \bar{A} \cap X$$
$$= \bar{A}$$

\noindent
(b - d)
Skip.

\subsection*{Exercise 20}
(a)
Its boundary is itself, and it has no interior.

\noindent
(b)
Its interior is itself, and its boundary is the union of line $x = 0$ with the line $y = 0$ for $x > 0$.

\noindent
(c)
The interior is the region $x > 0$, and its boundary is the union of line $x = 0$ with the line $y = 0$ for $x < 0$.

\noindent
(d)
Since both the rational and irrational are dense in $\mathbb{R}$, then the boundary is the entire $\mathbb{R}^2$. It has no interior.

\noindent
(e - f)
Use Desmos to see.

\subsection*{Exercise 21}
Closure is idempotent, and complementation is an involution. Thus, applying either twice in a row yields nothing new. Hence, denoting closure by C and complementation by S, we only have two chains of operations:
$$SCSC \ldots$$
$$CSCS \ldots$$

We already have $A$. There are 13 more to find. Idk how to proceed now.

\section{Continuous Functions}

\subsection*{Exercise 1}
Assuming the standard topology, consider a basis element of the codomain, which is an open interval $V$. For each $y \in V$, we look at $V_y(\epsilon) \subseteq V$. Using the epsilon-delta property, there exists $U_x(\delta)$ around each preimage $x \in f^{-1}(\{ y \})$ such that $f(U_x(\delta)) \subseteq V_y(\epsilon)$.

Claim: $\bigcup U_x$ is equal to the preimage of $V$. Right inclusion: every $U_x$ is mapped inside some $V_y \subseteq V$, hence it's part of the preimage of $V$. Left inclusion: every $x \in f^{-1}(V)$ is contained in some $U_x$.

Since each $U_x$ is open, their union is open. Hence, the preimage of open intervals is open, hence epsilon-delta implies topological continuity.

\subsection*{Exercise 2}
I found the solution online. The constant function is continuous, but the constant is not a limit point of itself.

\subsection*{Exercise 3}
(a)
Suppose $i$ is identity function. If $f$ is continuous, then for every $U \in \mathcal{T}$, we have $f^{-1}(U) = U$ is open, hence $U \in \mathcal{T}'$, so $\mathcal{T}'$ is finer than $\mathcal{T}$. For the other direction, suppose $\mathcal{T}'$ is finer than $\mathcal{T}$. If $U \in \mathcal{T}$, then $f^{-1}(U) = U \in \mathcal{T}'$, hence the preimage is open, hence $f$ is continuous.

\noindent
(b)
Well, if $i$ is a homeomorphism, then its inverse function is continuous, hence $\mathcal{T}$ is finer than $\mathcal{T}'$. Thus, $\mathcal{T} = \mathcal{T}'$.

\subsection*{Exercise 4}
Restricting our attention to the first function is enough. We are basically appending a constant to the representation of each $x \in X$. It's obviously invertible, and we can find its inverse $f^{-1}(x \times y_0) = x$. Each point in $X \times Y$ has a basis neighborhood $U_x \times U_y$ where $U_x$ and $U_y$ are open in $X$ and $Y$ respectively. Its restriction to the subspace $X \times y_0$ leaves basis element of the form $U_x \times y_0$. Using our inverse, its preimage $f^{-1}(U_x \times y_0)$ equals $U_x$, which is open in $X$. Likewise, if $U_x$ is open, then $f(U_x) = U_x \times y_0$, which is open in $X \times y_0$. Hence, $f$ is a homeomorphism from $X$ to $X \times y_0$, which is a subset of $X \times Y$, hence $f$ is an imbedding. Similar procedure for the second function.

\subsection*{Exercise 5}
Too lazy. Use $f(x) = \frac{x - a}{b - a}$.

\subsection*{Exercise 6}
$$f(x) = \begin{cases}
    x & \text{if } x \in \mathbb{Q} \\
    0 & \text{otherwise}
\end{cases}$$
Continuous only at $x = 0$.

\subsection*{Exercise 7}
(a)
The limit can be translated roughly to:
$$x \in [a, a + \delta) \rightarrow f(x) \in (f(a) - \epsilon, f(a) + \epsilon)$$
with $\delta$ being found for every $\epsilon$. Using the same reasoning as in Exercise 1, the preimage of $V = (f(a) - \epsilon, f(a) + \epsilon)$ is the union of half-open neighborhoods $U_x = [a, a + \delta)$ for each $x \in f^{-1}(\{ y \})$ for each $y \in V$. Since each $U_x$ is open in $\mathbb{R}_l$, then their union is also open. Hence, $f$ is continuous.

\noindent
(b)
Too lazy.

\section{The Product Topology}

\subsection*{Exercise 1}
The original basis for the box topology used every open subset from each $X_\alpha$. This alternative basis only uses the basis from each $X_\alpha$.

Anyways, given an original basis element $B$ of the box product space, we can decompose as $B = \prod U_\alpha$. For every $x \in B$, we decompose as $x = \prod x_\alpha$. For each $x_\alpha$, we can find a basis neighborhood $B_{x_\alpha} \subseteq U_\alpha$. Hence, $\prod B_{x_\alpha}$ is a neighborhood of $x$ fully contained inside $\prod U_\alpha$. The converse (every alternative basis can be expressed as original basis) is much easier, since there's complete freedom in choosing $U_\alpha$. Hence, this alternative basis is indeed a basis for the box topology.

For the product topology, we examine the basis generated by the usual subbasis, which is done by taking finite intersections:
$$\pi^{-1}_{\beta_1} (U_{\beta_1}) \cap \cdots \cap \pi^{-1}_{\beta_n} (U_{\beta_n})$$
Some of the reverse projections will come from the same $X_{\beta_i}$. In those case, we know that $\pi^{-1}_{\beta_i}(U_{\beta_i 1}) \cap \pi^{-1}_{\beta_i}(U_{\beta_i 2}) = \pi^{-1}_{\beta_i}(U_{\beta_i 1} \cap U_{\beta_i 2})$. Hence, each basis of the product topology can be generated by taking the reverse projection of at most one set from each affected $X_\beta$. The rest of the reasoning is similar to box topology.

\subsection*{Exercises 2 - 3}
Too lazy.

\subsection*{Exercise 4}
Use something similar to an identity map.

\subsection*{Exercise 5}
If $f$ continuous, then each $f_\alpha$ is continuous. Reason: Example 2.

\subsection*{Exercise 6}
Suppose $(x_i)$ converges to $x$. Given a neighborhood $B_{\pi_\alpha(x)}$, i.e. a neighborhood around the $\alpha$-th coordinate of $x$, we can form a neighborhood $B_x$ which is the product of: $X_\beta$ for every $\beta \neq \alpha$, and $B_{\pi_\alpha (x)}$ when $\beta = \alpha$. This is a basis in the product topology, hence it's a valid neighborhood of $x$. Hence, it contains every $x_{i > n}$ for some $n \in \mathbb{N}$. For each $x_{i > n}$, their $\alpha$-th coordinate is inside $B_{\pi_\alpha (x)}$. Hence, $\pi_\alpha(x_i)$ converges to $\pi_\alpha(x)$.

Conversely, suppose each sequence $(\pi_\alpha (x_i))$ converges to $\pi_\alpha(x)$. For every neighborhood $B_x$, we can express it as a finite product of neighborhoods of each coordinate $\prod B_{\pi_\alpha (x)}$. Each of them contains every $\pi_\alpha (x_{i > n_\alpha})$ for some $n_\alpha$. Taking their max, it follows that $B_x$ contains every $x_{i > n}$ for some $n = \max \{ n_\alpha \}$. Since each $B_x$ eventually contains every $x_i$, hence $x_i$ converges to $x$.

For the box topology, the above converse would be false, because we extract the maximum from a finite amount of neighborhoods, which may not exist for an infinite/uncountable number of neighborhoods allowed by the box topology.

\subsection*{Exercise 7}
In the box topology, take the sequence $x = (1, 1, \ldots)$ and let $B_x = (0.9, 1.1) \times (0.9, 1.1) \times \cdots$. None of the $y \in \mathbb{R}^\infty$ is contained in $B_x$, hence $x$ is not a limit point of $\mathbb{R}^\infty$. Hence, the closure of $\mathbb{R}^\infty$ in the box topology is itself.

On the other hand, every neighborhood of $x$ in the product topology constrains only a finite number of coordinates. Hence, choose $y \in \mathbb{R}^\infty$ that satisfies those constraints, and let the rest be $0$. Hence, the closure of $\mathbb{R}^\infty$ in the product topology is the entire product space.

\section{The Metric Topology}

\subsection*{Exercise 1}
(a)
The shapes are basically diamonds and octahedrons and higher dimensional ones. Every basis element of the product topology of standard topologies (which is the usual topology on $\mathbb{R}^n$) is the product of intervals $(a_1, b_1) \times \cdots \times (a_n, b_n)$. Then we can simply find $B_{d'}(\vec{x}, \epsilon)$ where $\vec{x}$ is a point inside this interval product, and $\epsilon$ is the minimum distance from $\vec{x}$ to each interval boundary. Since basis contained in basis, then the metric topology induced by $d'$ is finer than the product topology. Conversely, given an $\epsilon$-ball (or may I say, a diamond) around $\vec{x}$ in the metric topology, we choose each interval to be $(x_i - \frac{\epsilon}{n}, x_i + \frac{\epsilon}{n})$, their product becomes a basis element of the product topology. The maximal distance is achieved when a point approaches each boundary of every dimension, where its distance from $\vec{x}$ approaches $\frac{\epsilon}{n} + \cdots + \frac{\epsilon}{n} = \epsilon$. Of course, it stays below this value, because the intervals are open, so no point can rest exactly on the boundary to produce this value. Hence, the product topology is finer than the metric topology. Hence, they are equal.

\noindent
(b)
Too hard for my dumbass.

\subsection*{Exercise 2}
I found the solution online. This space is homeomorphic to $\mathbb{R}_d \times \mathbb{R}$. The metric that induces $\mathbb{R}_d$ is
$$d(x, y) = \begin{cases}
    0 & \text{if } x = y \\
    1 & \text{if } x \neq y
\end{cases}$$
which produces singletons. The metric on $\mathbb{R}$ is the standard metric. We combine them to produce a metric on their product, using the Manhattan distance for example. The metric would be
$$d(x, y) = \begin{cases}
    |x_2 - y_2| & \text{if } x_1 = y_1 \\
    |x_2 - y_2| + 1 & \text{if } x_1 \neq y_1
\end{cases}$$
It is kinda obvious now. We can create intervals of radius $1$ or smaller on each vertical stripe, which does not flood over to the neighboring stripes. We can take the union of a bunch of such intervals to recreate any basis element of the dictionary topology. Conversely, the balls generated by the metric are either intervals within a vertical stripe, or rectangles that are infinitely horizontal (but limited vertically). Clearly, we can take the union of intervals from the dictionary to recreate them.

\subsection*{Exercise 3}
(a)
Given a bounded basis element $V \subseteq \mathbb{R}$, which is an interval, we look at each $x_1 \times x_2 \in d^{-1}(\{ y \})$ for each $y \in V$. Our idea is to find a neighborhood $B_{x_1 \times x_2}$ that maps into $V$.

Consider a basis neighborhood around $x_1 \times x_2$, which can be decomposed as the product $B_{x_1 \times x_2} = B_\epsilon (x_1) \times B_\epsilon (x_2)$ where each $B_\epsilon \subseteq X$. We need to find a small enough $\epsilon$. For every pair $x_1' \times x_2' \in B_{x_1 \times x_2}$, notice that, by triangle inequality, we have
$$d(x_1', x_2') \leq d(x_1', x_1) + d(x_1, x_2) + d(x_2, x_2') \leq d(x_1, x_2) + 2\epsilon$$
By symmetry, we also have a lower bound
$$d(x_1, x_1') + d(x_1', x_2') + d(x_2', x_2) \geq d(x_1, x_2)$$
$$\Rightarrow d(x_1', x_2') + 2\epsilon \geq d(x_1, x_2)$$
$$\Rightarrow d(x_1', x_2') \geq d(x_1, x_2) - 2\epsilon$$
Put together, it looks like:
$$d(x_1, x_2) - 2\epsilon \leq d(x_1', x_2') \leq d(x_1, x_2) + 2\epsilon$$
By choosing
$$2\epsilon = \min \{ |\inf V - d(x_1, x_2)|, |\sup V - d(x_1, x_2)| \}$$
we guarantee that
$$\inf V \leq d(x_1', x_2') \leq \sup V$$
$$\Rightarrow d(x_1', x_2') \in [\inf V, \sup V]$$
If we make $\epsilon$ smaller (such as dividing by 2), it will fit inside the open interval version:
$$d(x_1', x_2') \in (\inf V, \sup V) \subseteq V$$

Hence, for every point $x = x_1 \times x_2$ in the preimage of $V$, we can find a neighborhood that maps inside $V$. Taking the union of these neighborhoods produces the preimage of $V$, which is thus also open. Hence the distance function is continuous.

\noindent
(b)
Too lazy.

\subsection*{Exercise 4}
(a)
In the product topology, since each component is continuous, all three functions are continuous.

In the uniform topology, only $g$ and $h$ are continuous. For $f$, use the neighborhood $B_{0.5}(0)$. No matter what interval $(a, b) \subseteq \mathbb{R}$ is chosen, one of the bounds will get scaled past $0.5$ at some coordinate eventually.

In the box topology, none of them are continuous. Since it's finer than the uniform topology, $f$ is already eliminated. A previous example in the textbook already eliminates $g$. But the general reasoning is that an interval in $\mathbb{R}$ must accomodate every dimension in $\mathbb{R}^\omega$. So for $h$, take the neighborhood $(-1, 1) \times (-\frac{1}{4}, \frac{1}{4}) \times (-\frac{1}{9}, \frac{1}{9}) \times \cdots$. The bounds shrink too fast for the harmonic sequence to catch up.

\noindent
(b)
In the product topology, $(w_n)$ converges to $\vec{0}$, because each coordinate converges to $0$ individually. It does not converge in the uniform topology, nor the box topology.

$(x_n)$ converges in the product topology under the same reasoning as $(w_n)$. However, it also converges to $\vec{0}$ under the uniform topology. For every neighborhood $B_\epsilon (\vec{0})$, we simply go far enough such that $\frac{1}{n} < \epsilon$. All the following $x_{i > n}$ will be contained in the neighborhood. Still, it does not converge in the box topology.

$(y_n)$ converges in the product topology under the same reasoning as $(w_n)$. It also converges in the uniform topology under the same reasoning as $(x_n)$. It doesn't converge in the box topology. For example, take the neighborhood $(-1, 1) \times (-\frac{1}{4}, \frac{1}{4}) \times (-\frac{1}{9}, \frac{1}{9}) \times \cdots$ from part (a).

$(z_n)$ converges in the product topology under the same reasoning as $(w_n)$. It also converges in the uniform topology under the same reasoning as $(x_n)$. It converges to $\vec{0}$ in the box topology, because it can be reduced to a 2-dimensional problem, and as we know, these topologies behave the same in finite-dimensional space.

\subsection*{Exercise 5}
Suppose a sequence $(x_n) \in \mathbb{R}^\omega$ converges to $L \neq 0$. Choose $\epsilon = \frac{|L|}{2}$ and form the neighborhood $B_\epsilon (x)$, where $x = (x_n)$. Visually, it looks like a horizontal band between $\frac{L}{2}$ and $\frac{3L}{2}$ in $\mathbb{R}^\omega$. Since every sequence in $\mathbb{R}^\infty$ is eventually zero, none of them is contained in this neighborhood. Hence, the closure of $\mathbb{R}^\infty$ in $\mathbb{R}^\omega$ must contain sequences that converge to zero at least. In fact, if a sequence $(x_n) \in \mathbb{R}^\omega$ converges to zero, then for every neighborhood $B_\epsilon (x)$, we can find some element $(y_n) \in \mathbb{R}^\infty$ inside the neighborhood. Since $(x_n)$ converges to zero, we can find $n \in \mathbb{N}$ such that $|x_{i > n}| < \epsilon$. For $i \leq n$, we set $y_i = x_i$, and for $i > n$, we set $y_i = 0$. Indeed, $y_i$ is eventually zero, hence it's an element of $\mathbb{R}^\infty$, but it's also contained in the neighborhood $B_\epsilon(x)$. Hence, the closure of $\mathbb{R}^\infty$ in $\mathbb{R}^\omega$ is every sequence that converges to zero.

\end{document}
